\documentclass[11pt]{article}

\usepackage[margin=1in]{geometry}
\usepackage{amsmath,amssymb,amsthm,mathtools,bm}
\usepackage[T1]{fontenc}
\usepackage[utf8]{inputenc}
\usepackage{lmodern}
\usepackage{microtype}
\usepackage[colorlinks=true,linkcolor=blue!50!black,urlcolor=blue!50!black]{hyperref}
\usepackage{xcolor}

\newtheorem{proposition}{Proposition}
\newtheorem{remark}{Remark}
\newtheorem{corollary}{Corollary}

\newcommand{\Complex}{\mathbb{C}}
\newcommand{\EE}{\mathbf{E}}
\newcommand{\hh}{\mathbf{h}}
\newcommand{\QQ}{\mathbf{Q}}
\newcommand{\x}{\mathbf{x}}
\newcommand{\I}{\mathbf{I}}
\newcommand{\Gh}{\widetilde{\mathbf{G}}}
\newcommand{\psin}{\boldsymbol{\psi}}
\newcommand{\norm}[1]{\left\lVert #1 \right\rVert}
\newcommand{\abs}[1]{\left\lvert #1 \right\rvert}
\newcommand{\dd}{\mathrm{d}}

\title{Direction 5: estimating $\QQ$ from channel feedback\\
\large Independent mathematical note}
\author{}
\date{2026-04-22}

\begin{document}
\maketitle

\section*{Setup}

The monograph defines the UE channel vector and exposure operator as
\begin{align}
  h_j
  &= \sum_{n:\,j(n)=j}
  \mathbf{C}_R(\hat{\mathbf{k}}_n)^H \psin_n\,
  e^{-ik_0\hat{\mathbf{k}}_n\cdot \mathbf{r}_{\mathrm{UE}}},
  \label{eq:h-note}
  \\
  \QQ
  &= \int_\Sigma \Gh(\mathbf{r})^H \Gh(\mathbf{r})\,\dd A,
  \qquad
  P_{\mathrm{abs}}(\x)=\x^H\QQ\x,
  \label{eq:Q-note}
\end{align}
where
\begin{equation}
  \Gh_j(\mathbf{r})
  = \sum_{n:\,j(n)=j}
  \sqrt{\frac{\sigma}{4\alpha_n}}\,
  \mathbf{F}_n(\mathbf{r})\psin_n\,
  e^{-ik_0\hat{\mathbf{k}}_n\cdot \mathbf{r}}.
\end{equation}
The ECBF precoder is
\begin{equation}
  \x^\star
  = \sqrt{P}\,
  \frac{(\lambda \QQ + \nu \I)^{-1}\hh^*}
  {\norm{(\lambda \QQ + \nu \I)^{-1}\hh^*}}.
\end{equation}

The question is not whether $\QQ$ is useful. It is. The question is whether
it can be inferred from standard channel feedback, and if not, what weaker
stability statements are enough for control.

\section*{1. Why $\hh$ does not identify $\QQ$}

\begin{proposition}[Generic non-identifiability]
The map from path-level data $\{(\hat{\mathbf{k}}_n,\psin_n)\}_n$ to the UE
channel $\hh$ is many-to-one. Consequently, $\QQ$ cannot in general be
recovered from $\hh$ alone.
\end{proposition}

\begin{proof}
Consider the simplest case with one path per BS element and fixed arrival
direction $\hat{\mathbf{k}}$. Then
\begin{equation}
  h_j = \mathbf{c}^H \psin_j\,e^{-i\phi_j},
  \qquad
  \mathbf{c} \equiv \mathbf{C}_R(\hat{\mathbf{k}}),
\end{equation}
while on the body
\begin{equation}
  \Gh_j(\mathbf{r})
  = \beta(\mathbf{r})\,\mathbf{F}(\mathbf{r})\psin_j\,e^{-i\phi_j(\mathbf{r})}.
\end{equation}
The receive map $\psin_j \mapsto \mathbf{c}^H\psin_j$ is scalar-valued, so
its nullspace is generically two-complex-dimensional. Pick any nonzero
$\delta_j \in \ker(\mathbf{c}^H)$ such that
$\mathbf{F}(\mathbf{r})\delta_j \not\equiv 0$ on a set of positive area.
Replacing $\psin_j$ by $\psin_j+\delta_j$ leaves $h_j$ unchanged because
$\mathbf{c}^H\delta_j=0$, but it changes the exposure channel
$\Gh_j(\mathbf{r})$, hence changes the Gram integral~\eqref{eq:Q-note}
generically.

So the same $\hh$ is compatible with infinitely many distinct exposure
operators. Algebraic inversion from $\hh$ to $\QQ$ is therefore impossible
without extra side information.
\end{proof}

\begin{remark}
This is not a dimension-counting heuristic. It is a constructive ambiguity:
the UE sees only a scalar projection of each path field, while the body
responds to the full Fresnel-filtered vector field over the entire surface.
\end{remark}

\section*{2. Why $\hh$ alone cannot certify small exposure alignment}

Define the monograph's alignment metric
\begin{equation}
  \rho(\hh,\QQ)
  \equiv
  \frac{\hh^T\QQ\,\hh^*}
  {\norm{\hh}^2\,\lambda_{\max}(\QQ)}
  \in [0,1].
\end{equation}

\begin{proposition}[No nontrivial $\hh$-only bound on $\rho$]
Fix any nonzero channel vector $\hh \in \Complex^M$ and any target
$\alpha \in [0,1]$. There exists a Hermitian PSD matrix $\QQ$ such that
\begin{equation}
  \rho(\hh,\QQ)=\alpha.
\end{equation}
\end{proposition}

\begin{proof}
Let $\mathbf{u}=\hh^*/\norm{\hh}$ and choose any unit vector
$\mathbf{v}\perp \mathbf{u}$. For any $\lambda>0$, define
\begin{equation}
  \QQ_\alpha
  = \lambda\bigl(\alpha\,\mathbf{u}\mathbf{u}^H
  + \mathbf{v}\mathbf{v}^H\bigr).
\end{equation}
Then $\QQ_\alpha$ is Hermitian PSD and
$\lambda_{\max}(\QQ_\alpha)=\lambda$. Also,
\begin{equation}
  \frac{\hh^T\QQ_\alpha\hh^*}{\norm{\hh}^2}
  = \mathbf{u}^H\QQ_\alpha\mathbf{u}
  = \alpha\lambda.
\end{equation}
Therefore $\rho(\hh,\QQ_\alpha)=\alpha$.
\end{proof}

\begin{remark}
This proposition is deliberately abstract: it ranges over all PSD matrices,
not only physically generated AEGIS operators. Its message is that
\emph{some additional structural prior is necessary}. Any claim of the form
``the channel alone implies ECBF is unnecessary'' must use more than
$\hh$ by itself.
\end{remark}

\section*{3. The right object for stale or approximate exposure control}

If exact recovery is impossible, the next question is whether an approximate
operator $\widehat{\QQ}$ is still good enough. The natural metric is the
spectral norm error
\begin{equation}
  \varepsilon_Q \equiv \norm{\QQ-\widehat{\QQ}}_2.
\end{equation}

\begin{proposition}[Uniform absorbed-power perturbation bound]
For any precoder $\x$ with $\norm{\x}^2 \le P$,
\begin{equation}
  \abs{\x^H\QQ\x - \x^H\widehat{\QQ}\x}
  \le P\,\varepsilon_Q.
\end{equation}
\end{proposition}

\begin{proof}
Let $\Delta=\QQ-\widehat{\QQ}$. Then
\begin{equation}
  \abs{\x^H\Delta\x}
  \le \norm{\Delta}_2\,\norm{\x}^2
  \le P\,\varepsilon_Q.
\end{equation}
\end{proof}

\begin{corollary}[Robust compliance margin]
If a design based on $\widehat{\QQ}$ enforces
\begin{equation}
  \x^H\widehat{\QQ}\x \le P_{\mathrm{abs}}^{\max} - P\,\varepsilon_Q,
\end{equation}
then the true absorbed power satisfies
\begin{equation}
  \x^H\QQ\x \le P_{\mathrm{abs}}^{\max}.
\end{equation}
\end{corollary}

\begin{proposition}[Worst-case exposure drift]
\begin{equation}
  \abs{\lambda_{\max}(\QQ)-\lambda_{\max}(\widehat{\QQ})}
  \le \varepsilon_Q.
\end{equation}
\end{proposition}

\begin{proof}
This is Weyl's inequality for Hermitian matrices.
\end{proof}

\begin{proposition}[Alignment drift]
Let
\begin{equation}
  \mu \equiv
  \min\!\bigl\{\lambda_{\max}(\QQ),\lambda_{\max}(\widehat{\QQ})\bigr\}
  > 0.
\end{equation}
Then
\begin{equation}
  \abs{\rho(\hh,\QQ)-\rho(\hh,\widehat{\QQ})}
  \le \frac{2\varepsilon_Q}{\mu}.
\end{equation}
\end{proposition}

\begin{proof}
Let $\mathbf{u}=\hh^*/\norm{\hh}$, and define
\begin{equation}
  a=\mathbf{u}^H\QQ\mathbf{u},
  \qquad
  \widehat{a}=\mathbf{u}^H\widehat{\QQ}\mathbf{u},
  \qquad
  b=\lambda_{\max}(\QQ),
  \qquad
  \widehat{b}=\lambda_{\max}(\widehat{\QQ}).
\end{equation}
Then $\rho(\hh,\QQ)=a/b$ and $\rho(\hh,\widehat{\QQ})=\widehat{a}/\widehat{b}$.
Moreover,
\begin{equation}
  \abs{a-\widehat{a}} \le \varepsilon_Q,
  \qquad
  \abs{b-\widehat{b}} \le \varepsilon_Q.
\end{equation}
Using $0\le \widehat{a}\le \widehat{b}$,
\begin{align}
  \abs{\frac{a}{b}-\frac{\widehat{a}}{\widehat{b}}}
  &\le \frac{\abs{a-\widehat{a}}}{b}
  + \widehat{a}\,\abs{\frac{1}{b}-\frac{1}{\widehat{b}}}
  \\
  &\le \frac{\varepsilon_Q}{\mu}
  + \widehat{b}\,\frac{\varepsilon_Q}{b\widehat{b}}
  \le \frac{2\varepsilon_Q}{\mu}.
\end{align}
\end{proof}

These three bounds already say something useful. If route A gives a stale
ray-traced operator $\widehat{\QQ}$, then the only quantity that needs
tracking for safety is not the full matrix error entrywise, but the induced
operator error $\varepsilon_Q$.

\section*{4. Pose drift as an operator perturbation}

The clean way to study orientation or posture staleness is to elevate the
body-surface channel to an operator. Let $\xi$ denote body pose
(position plus orientation, or just orientation if position is fixed), and
let
\begin{equation}
  \mathcal{T}_\xi:\Complex^M \to L^2(\Sigma;\Complex^3),
  \qquad
  (\mathcal{T}_\xi\x)(\mathbf{r})=\Gh_\xi(\mathbf{r})\x.
\end{equation}
Then
\begin{equation}
  \QQ(\xi)=\mathcal{T}_\xi^*\mathcal{T}_\xi.
\end{equation}

\begin{proposition}[Pose perturbation bound]
For any two poses $\xi,\xi_0$,
\begin{equation}
  \norm{\QQ(\xi)-\QQ(\xi_0)}_2
  \le
  \bigl(\norm{\mathcal{T}_\xi}_2+\norm{\mathcal{T}_{\xi_0}}_2\bigr)
  \norm{\mathcal{T}_\xi-\mathcal{T}_{\xi_0}}_2.
\end{equation}
\end{proposition}

\begin{proof}
Use
\begin{equation}
  \mathcal{T}_\xi^*\mathcal{T}_\xi
  -\mathcal{T}_{\xi_0}^*\mathcal{T}_{\xi_0}
  =
  (\mathcal{T}_\xi-\mathcal{T}_{\xi_0})^*\mathcal{T}_\xi
  + \mathcal{T}_{\xi_0}^*(\mathcal{T}_\xi-\mathcal{T}_{\xi_0}),
\end{equation}
then apply the triangle inequality and submultiplicativity of the operator
norm.
\end{proof}

\begin{corollary}[Lipschitz drift under smooth motion]
Assume
\begin{equation}
  \norm{\mathcal{T}_\xi-\mathcal{T}_{\xi_0}}_2
  \le L_T\,\norm{\xi-\xi_0}
\end{equation}
locally, and
\begin{equation}
  \sup_{\zeta \text{ near } \xi_0}\norm{\mathcal{T}_\zeta}_2 \le B.
\end{equation}
Then
\begin{equation}
  \norm{\QQ(\xi)-\QQ(\xi_0)}_2
  \le 2BL_T\,\norm{\xi-\xi_0}.
\end{equation}
\end{corollary}

\begin{remark}
This is the precise mathematical form of the intuitive statement
``if the body-surface channel changes smoothly with pose, then $\QQ$ also
changes smoothly with pose''. It reduces the staleness problem to estimating
one Lipschitz constant.
\end{remark}

\begin{corollary}[Safe refresh interval]
Suppose body pose evolves with speed $\norm{\dot{\xi}(t)}\le v_{\max}$ and
that $\QQ(\xi)$ is locally Lipschitz with constant $L_Q$. Over an update
interval $\Delta t$,
\begin{equation}
  \norm{\QQ(t+\Delta t)-\QQ(t)}_2 \le L_Q v_{\max}\Delta t.
\end{equation}
Hence any precoder with power budget $P$ remains conservatively safe if it
is designed against a margin
\begin{equation}
  P\,L_Q v_{\max}\Delta t.
\end{equation}
Equivalently, to guarantee an operator error budget $\varepsilon_Q$ it is
enough to refresh at
\begin{equation}
  \Delta t \le \frac{\varepsilon_Q}{L_Q v_{\max}}.
\end{equation}
\end{corollary}

For pure body rotation by angle $\theta$, this becomes
\begin{equation}
  \norm{\QQ(\theta)-\QQ(0)}_2 \le L_\theta \abs{\theta},
\end{equation}
so the absorbed-power error under any $\norm{\x}^2\le P$ obeys
\begin{equation}
  \abs{P_{\mathrm{abs}}(\x;\theta)-P_{\mathrm{abs}}(\x;0)}
  \le P\,L_\theta \abs{\theta}.
\end{equation}
This is exactly the kind of bound the markdown note asked for.

\section*{5. What seems most tractable next}

The mathematically clean path is not to recover $\QQ$ from $\hh$ alone.
It is to quantify how fast $\QQ$ moves on the pose manifold and then design
ECBF with a robustness margin.

That suggests the following AEGIS program:
\begin{enumerate}
  \item Fix a scene and body mesh, then compute $\QQ(\xi)$ over a grid of
  small pose perturbations using the existing coherent pipeline.
  \item Estimate $L_Q$ empirically from finite differences,
  \[
    L_Q \approx \max_{\xi,\xi'} \frac{\norm{\QQ(\xi)-\QQ(\xi')}_2}{\norm{\xi-\xi'}}.
  \]
  \item Compare the implied refresh time
  $\Delta t_{\max}=\varepsilon_Q/(L_Q v_{\max})$ against the radio channel
  coherence time.
  \item If $\Delta t_{\max}\gg T_c$, then a decoupled architecture becomes
  plausible: update $\hh$ every slot, update $\QQ$ every pose epoch.
\end{enumerate}

\section*{Bottom line}

My current view is:
\begin{itemize}
  \item Exact recovery of $\QQ$ from standard channel feedback is generically
  impossible.
  \item A nontrivial $\hh$-only bound on $\rho$ is impossible without extra
  priors on the admissible class of exposure operators.
  \item The right research question is therefore not inversion, but
  perturbation: how small can we make $\norm{\QQ-\widehat{\QQ}}_2$ under
  stale geometry or pose?
  \item Once that norm is controlled, whole-body absorbed power, worst-case
  eigenvalue, and alignment all admit direct quantitative bounds.
\end{itemize}

\end{document}
